\section{Conclusiones y limitaciones previstas}

La propuesta plantea caracterizar el uso de herramientas digitales entre estudiantes universitarios de Bogotá a partir de un enfoque cuantitativo, alcance descriptivo y diseño no experimental transversal. La encuesta estructurada es congruente con la necesidad de resumir tipos de herramientas, frecuencia, finalidades, utilidad percibida y dificultades reportadas. Esta elección permite una primera aproximación organizada al fenómeno, pero no permite probar relaciones causales ni medir directamente el aprendizaje.

La población y la muestra deben precisarse antes de la aplicación. Para sostener inferencias a una población definida se requiere un marco muestral y un procedimiento probabilístico compatible con los objetivos y recursos. Si solo es posible reclutar participantes voluntarios, el estudio deberá declarar que la muestra es no probabilística y limitar sus conclusiones a las respuestas obtenidas. El tamaño de la muestra no corrige por sí solo un sesgo de selección \parencite{argibay2009muestra}.

La revisión contextual muestra que la conectividad es relevante, pero los indicadores agregados de hogares no sustituyen la medición de experiencias estudiantiles. Por ello, el proyecto separa explícitamente el dato contextual del DANE de la información que deberá producir la encuesta \parencite{dane2025entic}. Del mismo modo, las recomendaciones sobre formación digital o apoyo institucional solo se formularán después de analizar los datos reales.

\subsection{Limitaciones previsibles}

Primero, el autoinforme puede introducir errores de recuerdo y respuestas socialmente deseables. Segundo, la encuesta en línea podría excluir a personas con dificultades de conectividad o acceso, por lo que conviene ofrecer una alternativa accesible si es viable. Tercero, la heterogeneidad entre instituciones y programas puede dificultar las comparaciones. Cuarto, el diseño transversal representa un periodo puntual y no permite identificar cambios ni establecer precedencia temporal. Quinto, el marco muestral y la institución participante todavía no están definidos.

\subsection{Acciones antes de la entrega final o ejecución}

\begin{enumerate}
    \item Completar los datos de portada y confirmar el nombre de la asignatura y del docente.
    \item Definir si se estudiará una institución específica o varias instituciones de Bogotá.
    \item Consultar directamente las lecturas obligatorias y completar los capítulos señalados por el curso, verificando las páginas y los conceptos que se citen.
    \item Obtener la autorización correspondiente y revisar el consentimiento informado.
    \item Revisar el cuestionario con el docente o un experto y efectuar una prueba piloto.
    \item Ajustar el plan de muestreo y análisis a los recursos reales disponibles.
\end{enumerate}

\subsection{Conclusión general}

El valor de este proyecto reside en producir una descripción transparente y metodológicamente delimitada del uso académico de herramientas digitales. Su calidad dependerá de la coherencia entre pregunta, objetivos, variables, muestra e instrumento, así como de la honestidad al informar las limitaciones. No se presentan resultados de campo porque la investigación aún no se ha ejecutado.

\subsection{Síntesis de la propuesta}


\subsection{Anexo breve: ejemplo de preguntas para el cuestionario}

\begin{table}[H]
    \caption{Ejemplo de preguntas para el cuestionario}
    \label{tab:preguntas-cuestionario}

    \small
    \renewcommand{\arraystretch}{1.3}

    \begin{tabularx}{\textwidth}{
        >{\raggedright\arraybackslash}p{0.22\textwidth}
        >{\raggedright\arraybackslash}X
    }
        \toprule
        \textbf{Dimensión} &
        \textbf{Pregunta propuesta} \\
        \midrule

        Herramientas &
        ¿Cuáles de estas herramientas utiliza para estudiar? (Marque todas las que correspondan). \\

        Frecuencia &
        En las últimas cuatro semanas, ¿con qué frecuencia utilizó herramientas digitales para actividades académicas? \\

        Finalidad &
        ¿Para qué actividades las utiliza principalmente? \\

        Utilidad &
        En una escala de 1 a 5, ¿qué tan útiles le resultan para completar sus tareas académicas? \\

        Dificultades &
        ¿Qué dificultades encuentra al utilizarlas? (Puede seleccionar varias). \\

        \bottomrule
    \end{tabularx}

    \vspace{2mm}
    \begin{minipage}{\textwidth}
        \footnotesize
        \textit{Nota.} Estas preguntas son ejemplos iniciales, no un instrumento validado. Deben revisarse y pilotearse antes de su uso.
    \end{minipage}
\end{table}
